%%%%%%%%%%%%%%%%%%%%%%%%%%%%%%%%%%%%%%%%%%%%%%%%%%%%%%%%%%%%%%%%%%%%%%%%
% Plantilla TFG/TFM
% Universidad de Murcia. Facultad de Informática
% Realizado por: José Manuel Requena Plens
% Modificado: Pablo José Rocamora Zamora
% Contacto: pablojoserocamora@gmail.com
%%%%%%%%%%%%%%%%%%%%%%%%%%%%%%%%%%%%%%%%%%%%%%%%%%%%%%%%%%%%%%%%%%%%%%%%


\chapter*{Extended Abstract}

Over the last few years, Large Language Models (LLMs) have changed the way information is produced, accessed, and shared in ways that would have seemed exaggerated just a decade ago. The fluency and apparent coherence of these systems is genuinely impressive, but it comes with a problem that becomes critical in certain applications: they hallucinate. They generate text that sounds authoritative and well-structured while being factually wrong, and they do so confidently. In a media landscape where misinformation already spreads faster than any manual verification effort can contain, handing a tool with these properties to newsrooms or fact-checkers without some kind of architectural safeguard is not a reasonable option.

This project starts from that premise. The goal was to design and evaluate a news verification system based on Retrieval-Augmented Generation (RAG), an architecture that tries to address the two core limitations of standalone language models. The first limitation is the static nature of their knowledge: once training ends, the model's understanding of the world is frozen, and no amount of prompting can make it know about events that happened after its training cutoff. The second is hallucination itself, the tendency to produce plausible but invented content when the model does not actually have reliable information to draw on. RAG addresses both by combining the model's language capabilities with a queryable external knowledge base, in this case a corpus of nearly half a million news articles indexed in Elasticsearch. The model does not answer from memory; it answers from retrieved evidence.

The project was structured around five specific objectives: deploying a retrieval infrastructure over the news corpus, comparing lexical and semantic search strategies, integrating generators of different scales (a local Small Language Model and two large cloud-based LLMs), establishing an automated evaluation framework using a 120-billion parameter judge model, and validating the final architecture on a curated 220-question benchmark. The overarching research question was whether RAG actually reduces hallucinations in a fact-checking setting, and how much the choice of retrieval method and generator model shape the final result.

The technical infrastructure was built on two remote university servers connected via SSH tunnels. The first handled storage and Elasticsearch operations; the second, equipped with an NVIDIA RTX 4090 GPU, was dedicated to computationally demanding tasks such as generating dense vector embeddings and running local model inference. All development was done in Python 3.12 within a Conda environment, using Jupyter Notebooks as the execution interface to document each phase iteratively. The choice of Elasticsearch (version 8.12.1) as the search engine was deliberate: unlike pure vector databases such as Pinecone or Milvus, Elasticsearch supports both BM25 keyword search and kNN vector search natively within the same platform, which made it ideal for the comparative experiment at the core of this project.

The raw dataset contained 804,083 news records, but not all of them were usable. An initial data audit showed that meaningful content resided almost exclusively in JSON-LD objects (around 62.5\% of the sample), and that the top-level \texttt{articleBody} field was systematically empty across the entire dataset. This shaped the extraction logic significantly: the system prioritized the nested \texttt{articleBody} field within the JSON-LD structure, falling back to \texttt{description} for multimedia-heavy records where the main body was absent. A cleaning function handled HTML entity decoding and whitespace normalization before indexing. Records with no textual content whatsoever (30,860 in total) were discarded, leaving a final index of 490,956 documents. The Elasticsearch index was configured with a Spanish language analyzer that performs lowercasing, stopword removal, and morphological stemming, optimizing it for the kind of lexical queries BM25 depends on. The corpus spans more than fifteen major outlets (La Vanguardia, El Universal, El Confidencial, BBC News, and others), which provides enough editorial diversity to reduce the risk of systematic bias in the retrieved evidence.

Two retrieval strategies were implemented and compared throughout the project. BM25 served as the lexical baseline: it ranks documents by term frequency and inverse document frequency, making it reliable for queries that contain proper nouns or specific dates, but less useful when the user rephrases a concept or uses synonyms. The semantic retrieval strategy used the BAAI/bge-m3 embedding model to represent both documents and queries as dense vectors in a 1,024-dimensional space, with kNN search using cosine similarity as the distance metric. This approach can locate conceptually relevant documents even when there is no lexical overlap between the query and the text.

One architectural decision worth explaining is the choice not to use chunking. Most RAG systems split documents into smaller fragments to fit within the language model's context window. Here, complete articles were indexed and injected instead. The reasoning was straightforward: news articles have a self-contained narrative structure, and splitting them risks separating entities or temporal references that belong together. The Phi-3 model's 4,096-token context window was sufficient to accommodate full articles in most cases, and evaluating documents in their entirety gives the retrieval engine more signal to work with when computing relevance scores. The Top-k hyperparameter was set to 2 after initial tests showed that a single retrieved document was sometimes insufficient for complex or ambiguous queries.

Three language models were evaluated as generators. Phi-3-mini-4k-instruct ran locally on the GPU node, making it the most computationally economical option. Llama 3.3 70B Versatile and Qwen 3 32B were accessed via the Groq API. All three were configured with temperature set to zero to make outputs deterministic and reproducible. Two distinct system prompts governed their behavior depending on the evaluation mode. In zero-shot mode (no external context), models received a minimal instruction to answer briefly based on their internal knowledge. In RAG mode, a strict fact-checking prompt was used, instructing the model to answer exclusively from the injected context and to respond with a fixed abstention phrase ("No tengo información suficiente en mis archivos") when the retrieved documents did not contain enough information to answer the query.

Evaluating 220 model responses manually would have been slow and hard to scale. The solution was a judge model: a 120-billion parameter LLM configured through LangChain to compare each generated response against the expected answer and return a binary verdict (CORRECT or INCORRECT). The judge was instructed to evaluate semantic equivalence rather than exact wording, and to base its verdict solely on the comparison between the two texts provided (not its own external knowledge). A human review pass validated the judge's verdicts after each batch, correcting the small number of cases where the automated evaluation was clearly off. Across the 80-question pilot phase, with five configurations evaluated (400 individual verdicts in total), the concordance between the judge and the human reviewer reached 97.75\% (9 discrepancies out of 400 individual verdicts), which gave sufficient confidence in the automated evaluation methodology going forward.

The pilot phase (80 questions) established a first hierarchy between configurations. Phi-3 without any context scored 12.5\%. Llama 3.3 70B in zero-shot mode reached 32.5\%. Qwen 3 32B in the same conditions scored exactly the same (32.5\%), which was notable: a 32-billion parameter model matched a 70-billion parameter model on general factual recall, suggesting that raw scale is not the determining factor once models reach a certain size. The RAG configuration with BM25 and Phi-3 brought the score up to 53.8%, and the semantic RAG with BGE-M3 and Phi-3 reached 92.5% on this initial subset. The pattern was already clear, though the full dataset would later moderate these figures.

The complete benchmark comprised 220 questions split into three categories designed to test different aspects of the system. Factual questions (143 total) asked for specific data points: dates, figures, names. Inference questions (37 total) required connecting information from different parts of a document or drawing a conclusion that was implied rather than stated explicitly. Trap questions (40 total) were the most interesting from a system reliability standpoint: they asked about events or facts not present in the corpus, meaning the correct response was always abstention. Getting these right requires the system to recognize the limits of its own knowledge base, which is a much harder problem than simply extracting information that is there.

Final results on the full 220-question dataset were as follows: Phi-3 zero-shot scored 16.4% (36/220), Llama 3.3 zero-shot 29.5% (65/220), Qwen 3 zero-shot 24.5% (54/220), RAG with BM25 and Phi-3 45.0% (99/220), RAG with BGE-M3 and Phi-3 70.9% (156/220), and the optimal system combining Llama 3.3 with BGE-M3 75.9% (167/220). The apparent gap between the BGE-M3 and Phi-3 result on the pilot (92.5%) and on the full dataset (70.9%) does not reflect a degradation in the system but a compositional shift in the benchmark. Two factors explain it: the pilot set contained a higher share of trap questions (22.5% of the first 80, versus 15.7% of the remaining 140), and trap questions are the category where the RAG system performs best due to its abstention instruction; additionally, the topics covered in the first 80 questions were high-density media events well-represented in the corpus (Trump tariffs, the Valencia floods, Euro 2024, Taylor Swift), whereas later questions targeted more niche subjects --- the Genoa maritime tunnel, Alzheimer drug trials, Neuralink, the Athena space probe --- where relevant documents were sparser and harder to retrieve. The jump from zero-shot to RAG (even with the simplest lexical retrieval) was larger than any improvement achievable by switching between generators, which says something important about where the real bottleneck lies in this kind of system.

Beyond aggregate numbers, the error analysis revealed something arguably more meaningful than the accuracy gap. Among the errors made by the Phi-3 semantic RAG, 32.3% were hallucinations: cases where the model produced an incorrect answer rather than staying silent. Llama 3.3, operating within the exact same RAG setup, brought that figure down to 2.2%. This is not just a quantitative difference; it reflects a qualitative shift in what the system does when it fails. Phi-3, being a smaller model with less alignment training, tends to fall back on its parametric memory when the system prompt conflicts with what it "knows", generating a confident but wrong answer. Llama 3.3, which already showed strong abstention behavior in zero-shot conditions (52.5% correct on trap questions, compared to under 23\% for the other two models), carries that tendency into the RAG setup and turns most retrieval failures into safe abstentions instead of dangerous hallucinations. On trap questions specifically, the optimal system reached 95% accuracy, meaning it correctly identified the absence of information in the corpus in 38 out of 40 cases.

A handful of case studies were included to illustrate both the strengths and the failure modes of the architecture. A zero-shot Llama 3.3 70B confidently stated that US tariffs on China reached 25\% when the correct figure was 145\%; the semantic RAG extracted the right number directly from the retrieved article. In a trap question about a German football coach, the zero-shot model invented a specific name while the RAG correctly abstained. A third case illustrated the reverse anomaly: when asked whether the European version of Meta AI had been trained on European user data, Phi-3 in zero-shot mode gave the correct answer (no), while Llama 3.3 70B --- despite its larger scale --- incorrectly stated that it had. This shows that aggregate accuracy differences between models do not prevent smaller models from outperforming larger ones on individual queries. On the other hand, the system also failed in two instructive ways: when the retrieval engine did not find the document containing the answer, the model was forced to abstain even though it might have known the fact from training (the question about the actor in a 1953 film mentioned in an Oscars article); and when a retrieved article contained a typographic error (a news piece incorrectly stating that Elon Musk founded Neuralink in 2026), the RAG reproduced that error faithfully, since it had no way to distinguish between a factual claim and a typo in the source.

The optimal architecture selected was the combination of BGE-M3 for retrieval and Llama 3.3 70B as the generator, not simply because they were the best-performing components individually, but because of how they complement each other. The selection of Llama 3.3 70B over the other generators was grounded in a specific criterion: its zero-shot performance on trap questions (52.5%), far above Phi-3 and Qwen 3 32B (both below 23%). In a RAG system, the generator operates under a strict abstention instruction; a model that already tends to stay silent when uncertain will follow that instruction more reliably than one that tends to fabricate. This makes trap question behavior in zero-shot mode a better predictor of RAG reliability than overall accuracy. Semantic retrieval minimizes the cases where relevant evidence goes unfound; Llama 3.3's alignment toward abstention ensures that when retrieval does fail, the system stays silent rather than fabricating an answer. In a fact-checking context, that distinction matters enormously.

Looking ahead, three directions seem most worth pursuing. Extending the system to an agentic RAG architecture with real-time web access would remove the temporal constraint of the static corpus and make it applicable to breaking news. Applying fine-tuning to a lightweight local model like Phi-3 to teach it Llama 3.3's abstention behavior would make high-reliability fact-checking feasible without cloud API dependencies. And incorporating multimodal capabilities would open the door to verifying image-based misinformation, which represents a growing share of what spreads online.

One thing this project reinforced is that systems designed to assess what is true carry ethical weight that purely technical metrics do not capture. A RAG system assumes that what has been published is accurate, which is not always the case. Editorial bias in the corpus can be reproduced and amplified, especially in inference tasks where the model connects dots across a text written from a particular perspective. This system was built as a tool to assist human fact-checkers, not to replace them, and that framing is not just a disclaimer: it reflects a genuine design constraint. Keeping a person in the loop is not a limitation of the current version that should be engineered away; it is a feature.
